\chapter{Unit 2: Fundamentals of A.C. Circuits}

\section{Section A: Multiple Choice Question Practice (60 Curated Problems)}

\noindent\textbf{Q1. An alternating current is one whose magnitude and direction:} \hfill [Topic: alternating current and voltage | Easy]
\begin{enumerate}[label=(\alph*),topsep=2pt,itemsep=1pt,partopsep=0pt,parsep=0pt]
    \item {Remain constant with time}
    \item {Change periodically with time}
    \item {Only change in magnitude but not direction}
    \item {Only change in direction but not magnitude}
\end{enumerate}
\vspace{2mm}

\noindent\textbf{Q2. The standard waveform assumed for an AC supply in analysis is:} \hfill [Topic: alternating current and voltage | Easy]
\begin{enumerate}[label=(\alph*),topsep=2pt,itemsep=1pt,partopsep=0pt,parsep=0pt]
    \item {Triangular}
    \item {Sinusoidal}
    \item {Sawtooth}
    \item {Square}
\end{enumerate}
\vspace{2mm}

\noindent\textbf{Q3. In AC circuit notation, lowercase letters such as i and v generally represent:} \hfill [Topic: concept of notations ( i, v, I, V) | Easy]
\begin{enumerate}[label=(\alph*),topsep=2pt,itemsep=1pt,partopsep=0pt,parsep=0pt]
    \item {Instantaneous values}
    \item {Maximum values}
    \item {RMS values}
    \item {Average values}
\end{enumerate}
\vspace{2mm}

\noindent\textbf{Q4. Uppercase symbols I and V in AC analysis most commonly represent:} \hfill [Topic: concept of notations ( i, v, I, V) | Easy]
\begin{enumerate}[label=(\alph*),topsep=2pt,itemsep=1pt,partopsep=0pt,parsep=0pt]
    \item {Angular frequency}
    \item {Peak-to-peak values}
    \item {RMS (effective) values}
    \item {Instantaneous values}
\end{enumerate}
\vspace{2mm}

\noindent\textbf{Q5. The amplitude of an AC waveform refers to its:} \hfill [Topic: definitions of amplitude, phase, phase difference | Easy]
\begin{enumerate}[label=(\alph*),topsep=2pt,itemsep=1pt,partopsep=0pt,parsep=0pt]
    \item {Average value}
    \item {RMS value}
    \item {Time period}
    \item {Maximum (peak) value}
\end{enumerate}
\vspace{2mm}

\noindent\textbf{Q6. Phase difference between two AC quantities is usually measured in:} \hfill [Topic: definitions of amplitude, phase, phase difference | Easy]
\begin{enumerate}[label=(\alph*),topsep=2pt,itemsep=1pt,partopsep=0pt,parsep=0pt]
    \item {Amperes}
    \item {Degrees or radians}
    \item {Volts}
    \item {Seconds only}
\end{enumerate}
\vspace{2mm}

\noindent\textbf{Q7. Two AC waveforms that reach their peaks at the same instant are said to be:} \hfill [Topic: definitions of amplitude, phase, phase difference | Easy]
\begin{enumerate}[label=(\alph*),topsep=2pt,itemsep=1pt,partopsep=0pt,parsep=0pt]
    \item {In phase}
    \item {Out of phase by 18 0 $^\circ$}
    \item {In quadrature}
    \item {Out of phase by 9 0 $^\circ$}
\end{enumerate}
\vspace{2mm}

\noindent\textbf{Q8. For a sinusoidal waveform of peak value V m , the RMS value is:} \hfill [Topic: RMS value and average value of an AC signal | Easy]
\begin{enumerate}[label=(\alph*),topsep=2pt,itemsep=1pt,partopsep=0pt,parsep=0pt]
    \item {2 V m}
    \item {2 V m}
    \item {2 V m}
    \item {$\pi$ 2 V m}
\end{enumerate}
\vspace{2mm}

\noindent\textbf{Q9. The average value of a sinusoidal waveform over a complete cycle is:} \hfill [Topic: RMS value and average value of an AC signal | Easy]
\begin{enumerate}[label=(\alph*),topsep=2pt,itemsep=1pt,partopsep=0pt,parsep=0pt]
    \item {0.637 V m}
    \item {Zero}
    \item {0.707 V m}
    \item {V m}
\end{enumerate}
\vspace{2mm}

\noindent\textbf{Q10. The average value of a sinusoidal waveform over a half cycle in terms of peak value V m is:} \hfill [Topic: RMS value and average value of an AC signal | Easy]
\begin{enumerate}[label=(\alph*),topsep=2pt,itemsep=1pt,partopsep=0pt,parsep=0pt]
    \item {0.637 V m}
    \item {0.5 V m}
    \item {0.318 V m}
    \item {0.707 V m}
\end{enumerate}
\vspace{2mm}

\noindent\textbf{Q11. The impedance of a circuit in complex (rectangular) form is written as:} \hfill [Topic: complex representation of impedance | Easy]
\begin{enumerate}[label=(\alph*),topsep=2pt,itemsep=1pt,partopsep=0pt,parsep=0pt]
    \item {Z = R + j X}
    \item {Z = j R}
    \item {Z = R $\times$ X}
    \item {Z = R - R}
\end{enumerate}
\vspace{2mm}

\noindent\textbf{Q12. In a purely inductive circuit, the inductive reactance is given by:} \hfill [Topic: complex representation of impedance | Easy]
\begin{enumerate}[label=(\alph*),topsep=2pt,itemsep=1pt,partopsep=0pt,parsep=0pt]
    \item {X L = $\omega$ C 1}
    \item {X L = $\omega$ L 1}
    \item {X L = $\omega$ C}
    \item {X L = $\omega$ L}
\end{enumerate}
\vspace{2mm}

\noindent\textbf{Q13. The capacitive reactance X C of a capacitor is expressed as:} \hfill [Topic: complex representation of impedance | Easy]
\begin{enumerate}[label=(\alph*),topsep=2pt,itemsep=1pt,partopsep=0pt,parsep=0pt]
    \item {$\omega$ L}
    \item {$\omega$ L 1}
    \item {$\omega$ C 1}
    \item {$\omega$ C}
\end{enumerate}
\vspace{2mm}

\noindent\textbf{Q14. In a series RL circuit, the current:} \hfill [Topic: steady state analysis of ac circuits consisting of RL, RC and RLC (series) | Easy]
\begin{enumerate}[label=(\alph*),topsep=2pt,itemsep=1pt,partopsep=0pt,parsep=0pt]
    \item {Leads the applied voltage}
    \item {Is zero}
    \item {Lags the applied voltage}
    \item {Is in phase with the voltage}
\end{enumerate}
\vspace{2mm}

\noindent\textbf{Q15. In a series RC circuit, the current:} \hfill [Topic: steady state analysis of ac circuits consisting of RL, RC and RLC (series) | Easy]
\begin{enumerate}[label=(\alph*),topsep=2pt,itemsep=1pt,partopsep=0pt,parsep=0pt]
    \item {Lags the applied voltage}
    \item {Is always maximum}
    \item {Is in phase with the voltage}
    \item {Leads the applied voltage}
\end{enumerate}
\vspace{2mm}

\noindent\textbf{Q16. The magnitude of impedance in a series RLC circuit is given by:} \hfill [Topic: steady state analysis of ac circuits consisting of RL, RC and RLC (series) | Easy]
\begin{enumerate}[label=(\alph*),topsep=2pt,itemsep=1pt,partopsep=0pt,parsep=0pt]
    \item {R 2 + ( X L - X C ) 2}
    \item {R 2 + X L 2 + X C 2}
    \item {R 2 + ( X L - X C ) 2}
    \item {R + X L + X C}
\end{enumerate}
\vspace{2mm}

\noindent\textbf{Q17. At resonance in a series RLC circuit, the net reactance is:} \hfill [Topic: resonance in series RLC circuit | Easy]
\begin{enumerate}[label=(\alph*),topsep=2pt,itemsep=1pt,partopsep=0pt,parsep=0pt]
    \item {Infinite}
    \item {Equal to R}
    \item {Maximum}
    \item {Zero ( X L = X C )}
\end{enumerate}
\vspace{2mm}

\noindent\textbf{Q18. The resonant frequency of a series RLC circuit is given by:} \hfill [Topic: resonance in series RLC circuit | Easy]
\begin{enumerate}[label=(\alph*),topsep=2pt,itemsep=1pt,partopsep=0pt,parsep=0pt]
    \item {f r = 2 $\pi$ L C 1}
    \item {f r = 2 $\pi$ L C 1}
    \item {f r = 2 $\pi$ L C}
    \item {f r = 2 $\pi$ L C}
\end{enumerate}
\vspace{2mm}

\noindent\textbf{Q19. The power factor of an AC circuit is defined as:} \hfill [Topic: power factor and power calculation in RL, RC and RLC circuits | Easy]
\begin{enumerate}[label=(\alph*),topsep=2pt,itemsep=1pt,partopsep=0pt,parsep=0pt]
    \item {sec $\phi$}
    \item {tan $\phi$}
    \item {sin $\phi$}
    \item {cos $\phi$}
\end{enumerate}
\vspace{2mm}

\noindent\textbf{Q20. In a balanced three-phase star (Y) connection, the relation between line voltage and phase voltage is:} \hfill [Topic: relations in line and phase voltages and currents in star and delta networks | Easy]
\begin{enumerate}[label=(\alph*),topsep=2pt,itemsep=1pt,partopsep=0pt,parsep=0pt]
    \item {V L = V p h}
    \item {V L = 3 V p h}
    \item {V L = 3 V p h}
    \item {V L = 3 V p h}
\end{enumerate}
\vspace{2mm}

\noindent\textbf{Q21. A sinusoidal voltage has a peak value of 200 V . What is its RMS value?} \hfill [Topic: RMS value and average value of an AC signal | Medium]
\begin{enumerate}[label=(\alph*),topsep=2pt,itemsep=1pt,partopsep=0pt,parsep=0pt]
    \item {127.3 V}
    \item {141.4 V}
    \item {100 V}
    \item {200 V}
\end{enumerate}
\vspace{2mm}

\noindent\textbf{Q22. The average value of a full sinusoidal AC current over one complete cycle is:} \hfill [Topic: RMS value and average value of an AC signal | Medium]
\begin{enumerate}[label=(\alph*),topsep=2pt,itemsep=1pt,partopsep=0pt,parsep=0pt]
    \item {Zero}
    \item {I m}
    \item {0.707 I m}
    \item {0.637 I m}
\end{enumerate}
\vspace{2mm}

\noindent\textbf{Q23. The form factor of a pure sinusoidal waveform is approximately:} \hfill [Topic: RMS value and average value of an AC signal | Medium]
\begin{enumerate}[label=(\alph*),topsep=2pt,itemsep=1pt,partopsep=0pt,parsep=0pt]
    \item {0.637}
    \item {0.707}
    \item {1.11}
    \item {1.414}
\end{enumerate}
\vspace{2mm}

\noindent\textbf{Q24. In a pure capacitor, the current leads the applied voltage by a phase angle of:} \hfill [Topic: phase difference | Medium]
\begin{enumerate}[label=(\alph*),topsep=2pt,itemsep=1pt,partopsep=0pt,parsep=0pt]
    \item {4 5 $^\circ$}
    \item {9 0 $^\circ$}
    \item {0 $^\circ$}
    \item {18 0 $^\circ$}
\end{enumerate}
\vspace{2mm}

\noindent\textbf{Q25. A current is expressed as i = 10 sin ( 314 t + 3 0 $^\circ$ ) A . What is the frequency of the supply?} \hfill [Topic: phase | Medium]
\begin{enumerate}[label=(\alph*),topsep=2pt,itemsep=1pt,partopsep=0pt,parsep=0pt]
    \item {100 Hz}
    \item {60 Hz}
    \item {50 Hz}
    \item {314 Hz}
\end{enumerate}
\vspace{2mm}

\noindent\textbf{Q26. A circuit has resistance R = 3 $\Omega$ and inductive reactance X L = 4 $\Omega$ . The magnitude of the impedance is:} \hfill [Topic: complex representation of impedance | Medium]
\begin{enumerate}[label=(\alph*),topsep=2pt,itemsep=1pt,partopsep=0pt,parsep=0pt]
    \item {5 $\Omega$}
    \item {7 $\Omega$}
    \item {1 $\Omega$}
    \item {25 $\Omega$}
\end{enumerate}
\vspace{2mm}

\noindent\textbf{Q27. The impedance of a series RC circuit is written in complex form as:} \hfill [Topic: complex representation of impedance | Medium]
\begin{enumerate}[label=(\alph*),topsep=2pt,itemsep=1pt,partopsep=0pt,parsep=0pt]
    \item {Z = - R + j X C}
    \item {Z = R + j X C}
    \item {Z = R + j X L}
    \item {Z = R - j X C}
\end{enumerate}
\vspace{2mm}

\noindent\textbf{Q28. In a series RL circuit, R = 10 $\Omega$ and X L = 10 $\Omega$ . The phase angle of the circuit is:} \hfill [Topic: steady state analysis of ac circuits consisting of RL, RC and RLC (series) | Medium]
\begin{enumerate}[label=(\alph*),topsep=2pt,itemsep=1pt,partopsep=0pt,parsep=0pt]
    \item {4 5 $^\circ$ leading}
    \item {6 0 $^\circ$ lagging}
    \item {4 5 $^\circ$ lagging}
    \item {3 0 $^\circ$ lagging}
\end{enumerate}
\vspace{2mm}

\noindent\textbf{Q29. A series RLC circuit has R = 6 $\Omega$ , X L = 12 $\Omega$ , and X C = 4 $\Omega$ . The net impedance magnitude is:} \hfill [Topic: steady state analysis of ac circuits consisting of RL, RC and RLC (series) | Medium]
\begin{enumerate}[label=(\alph*),topsep=2pt,itemsep=1pt,partopsep=0pt,parsep=0pt]
    \item {22 $\Omega$}
    \item {10 $\Omega$}
    \item {8 $\Omega$}
    \item {14 $\Omega$}
\end{enumerate}
\vspace{2mm}

\noindent\textbf{Q30. At resonance in a series RLC circuit, the impedance of the circuit is:} \hfill [Topic: resonance in series RLC circuit | Medium]
\begin{enumerate}[label=(\alph*),topsep=2pt,itemsep=1pt,partopsep=0pt,parsep=0pt]
    \item {Minimum and purely reactive}
    \item {Maximum and purely reactive}
    \item {Minimum and purely resistive}
    \item {Maximum and purely resistive}
\end{enumerate}
\vspace{2mm}

\noindent\textbf{Q31. The resonant frequency of a series RLC circuit with L = 2 H and C = 2 $\mu$ F is approximately:} \hfill [Topic: resonance in series RLC circuit | Medium]
\begin{enumerate}[label=(\alph*),topsep=2pt,itemsep=1pt,partopsep=0pt,parsep=0pt]
    \item {50 Hz}
    \item {159 Hz}
    \item {79.6 Hz}
    \item {500 Hz}
\end{enumerate}
\vspace{2mm}

\noindent\textbf{Q32. At resonance, the current in a series RLC circuit is in phase with the voltage because:} \hfill [Topic: resonance in series RLC circuit | Medium]
\begin{enumerate}[label=(\alph*),topsep=2pt,itemsep=1pt,partopsep=0pt,parsep=0pt]
    \item {The current becomes zero}
    \item {The circuit behaves as purely resistive}
    \item {The circuit behaves as purely capacitive}
    \item {The circuit behaves as purely inductive}
\end{enumerate}
\vspace{2mm}

\noindent\textbf{Q33. An RL circuit draws 5 A at 230 V with a power factor of 0.8 . The real power consumed is:} \hfill [Topic: power factor and power calculation in RL, RC and RLC circuits | Medium]
\begin{enumerate}[label=(\alph*),topsep=2pt,itemsep=1pt,partopsep=0pt,parsep=0pt]
    \item {1437 W}
    \item {690 W}
    \item {1150 W}
    \item {920 W}
\end{enumerate}
\vspace{2mm}

\noindent\textbf{Q34. The power factor of a circuit having impedance Z = 8 + j 6 $\Omega$ is:} \hfill [Topic: power factor and power calculation in RL, RC and RLC circuits | Medium]
\begin{enumerate}[label=(\alph*),topsep=2pt,itemsep=1pt,partopsep=0pt,parsep=0pt]
    \item {0.6 lagging}
    \item {0.8 lagging}
    \item {0.6 leading}
    \item {0.8 leading}
\end{enumerate}
\vspace{2mm}

\noindent\textbf{Q35. In a purely inductive AC circuit, the average power consumed over one cycle is:} \hfill [Topic: power factor and power calculation in RL, RC and RLC circuits | Medium]
\begin{enumerate}[label=(\alph*),topsep=2pt,itemsep=1pt,partopsep=0pt,parsep=0pt]
    \item {0.5 V I}
    \item {V I}
    \item {0.707 V I}
    \item {Zero}
\end{enumerate}
\vspace{2mm}

\noindent\textbf{Q36. In a balanced star-connected system, the relation between line voltage and phase voltage is:} \hfill [Topic: relations in line and phase voltages and currents in star and delta networks | Medium]
\begin{enumerate}[label=(\alph*),topsep=2pt,itemsep=1pt,partopsep=0pt,parsep=0pt]
    \item {V L = V p h}
    \item {V L = 3 V p h}
    \item {V L = 3 V p h}
    \item {V L = 3 V p h}
\end{enumerate}
\vspace{2mm}

\noindent\textbf{Q37. In a balanced delta-connected system, the relation between line current and phase current is:} \hfill [Topic: relations in line and phase voltages and currents in star and delta networks | Medium]
\begin{enumerate}[label=(\alph*),topsep=2pt,itemsep=1pt,partopsep=0pt,parsep=0pt]
    \item {I L = 3 I p h}
    \item {I L = 3 I p h}
    \item {I L = I p h}
    \item {I L = 3 I p h}
\end{enumerate}
\vspace{2mm}

\noindent\textbf{Q38. A balanced star-connected load has a line voltage of 400 V . The phase voltage is approximately:} \hfill [Topic: relations in line and phase voltages and currents in star and delta networks | Medium]
\begin{enumerate}[label=(\alph*),topsep=2pt,itemsep=1pt,partopsep=0pt,parsep=0pt]
    \item {692 V}
    \item {133 V}
    \item {231 V}
    \item {400 V}
\end{enumerate}
\vspace{2mm}

\noindent\textbf{Q39. In a three-phase balanced system, the phase difference between any two phase voltages is:} \hfill [Topic: three-phase circuits- numbering and interconnection (delta or mesh connection) of three phases | Medium]
\begin{enumerate}[label=(\alph*),topsep=2pt,itemsep=1pt,partopsep=0pt,parsep=0pt]
    \item {9 0 $^\circ$}
    \item {6 0 $^\circ$}
    \item {18 0 $^\circ$}
    \item {12 0 $^\circ$}
\end{enumerate}
\vspace{2mm}

\noindent\textbf{Q40. If an AC voltage completes 100 cycles in 2 seconds, its time period is:} \hfill [Topic: alternating current and voltage | Medium]
\begin{enumerate}[label=(\alph*),topsep=2pt,itemsep=1pt,partopsep=0pt,parsep=0pt]
    \item {0.5 s}
    \item {2 s}
    \item {50 s}
    \item {0.02 s}
\end{enumerate}
\vspace{2mm}

\noindent\textbf{Q41. A half-wave rectified sinusoidal current has a peak value of I m . What is the ratio of its RMS value to its average value?} \hfill [Topic: RMS value and average value of an AC signal | Hard]
\begin{enumerate}[label=(\alph*),topsep=2pt,itemsep=1pt,partopsep=0pt,parsep=0pt]
    \item {2 $\pi$}
    \item {$\pi$ 2}
    \item {2}
    \item {2 $\pi$}
\end{enumerate}
\vspace{2mm}

\noindent\textbf{Q42. A current waveform is given by i ( t ) = 10 + 5 sin ( $\omega$ t ) A. What is the RMS value of this current?} \hfill [Topic: RMS value and average value of an AC signal | Hard]
\begin{enumerate}[label=(\alph*),topsep=2pt,itemsep=1pt,partopsep=0pt,parsep=0pt]
    \item {112.5 A}
    \item {10.6 A}
    \item {125 A}
    \item {12.5 A}
\end{enumerate}
\vspace{2mm}

\noindent\textbf{Q43. A series RLC circuit has R = 10 $\Omega$ , L = 0.1 H, and C = 10 $\mu$ F. What is the quality factor Q at resonance?} \hfill [Topic: resonance in series RLC circuit | Hard]
\begin{enumerate}[label=(\alph*),topsep=2pt,itemsep=1pt,partopsep=0pt,parsep=0pt]
    \item {31.6}
    \item {1}
    \item {100}
    \item {10}
\end{enumerate}
\vspace{2mm}

\noindent\textbf{Q44. In a series RLC circuit at resonance, the voltage across the inductor is 50 V and across the resistor is 5 V. What is the quality factor Q ?} \hfill [Topic: resonance in series RLC circuit | Hard]
\begin{enumerate}[label=(\alph*),topsep=2pt,itemsep=1pt,partopsep=0pt,parsep=0pt]
    \item {5}
    \item {10}
    \item {50}
    \item {0.1}
\end{enumerate}
\vspace{2mm}

\noindent\textbf{Q45. An impedance is given as Z = 8 + j 6 $\Omega$ . If this is connected across a source V = 100$\angle$3 0 $^\circ$ V, what is the phase angle of the current?} \hfill [Topic: complex representation of impedance | Hard]
\begin{enumerate}[label=(\alph*),topsep=2pt,itemsep=1pt,partopsep=0pt,parsep=0pt]
    \item {36.8 7 $^\circ$}
    \item {66.8 7 $^\circ$}
    \item {- 36.8 7 $^\circ$}
    \item {- 6.8 7 $^\circ$}
\end{enumerate}
\vspace{2mm}

\noindent\textbf{Q46. A series RL circuit with R = 30 $\Omega$ and X L = 40 $\Omega$ is fed by 230 V, 50 Hz. What is the voltage across the inductor?} \hfill [Topic: steady state analysis of ac circuits consisting of RL, RC and RLC (series) | Hard]
\begin{enumerate}[label=(\alph*),topsep=2pt,itemsep=1pt,partopsep=0pt,parsep=0pt]
    \item {161 V}
    \item {230 V}
    \item {184 V}
    \item {138 V}
\end{enumerate}
\vspace{2mm}

\noindent\textbf{Q47. A single-phase load draws 10 A at 230 V with a power factor of 0.8 lagging. What is the reactive power consumed?} \hfill [Topic: power factor and power calculation in RL, RC and RLC circuits | Hard]
\begin{enumerate}[label=(\alph*),topsep=2pt,itemsep=1pt,partopsep=0pt,parsep=0pt]
    \item {920 VAR}
    \item {1840 VAR}
    \item {2300 VAR}
    \item {1380 VAR}
\end{enumerate}
\vspace{2mm}

\noindent\textbf{Q48. A balanced delta-connected load has a phase current of 10 A. If the same load impedances are reconnected in star across the same line voltage, what is the new line current?} \hfill [Topic: relations in line and phase voltages and currents in star and delta networks | Hard]
\begin{enumerate}[label=(\alph*),topsep=2pt,itemsep=1pt,partopsep=0pt,parsep=0pt]
    \item {3 10 3 A}
    \item {30 A}
    \item {10 3 A}
    \item {10 A}
\end{enumerate}
\vspace{2mm}

\noindent\textbf{Q49. In a balanced three-phase star-connected system, the line voltage is 400 V. What is the phase voltage?} \hfill [Topic: relations in line and phase voltages and currents in star and delta networks | Hard]
\begin{enumerate}[label=(\alph*),topsep=2pt,itemsep=1pt,partopsep=0pt,parsep=0pt]
    \item {400 V}
    \item {693 V}
    \item {566 V}
    \item {231 V}
\end{enumerate}
\vspace{2mm}

\noindent\textbf{Q50. A series RC circuit draws 5 A from a 100 V source and dissipates 300 W. What is the capacitive reactance X C ?} \hfill [Topic: power factor and power calculation in RL, RC and RLC circuits | Hard]
\begin{enumerate}[label=(\alph*),topsep=2pt,itemsep=1pt,partopsep=0pt,parsep=0pt]
    \item {20 $\Omega$}
    \item {8 $\Omega$}
    \item {16 $\Omega$}
    \item {12 $\Omega$}
\end{enumerate}
\vspace{2mm}

\noindent\textbf{Q51. A series RLC circuit has R = 5 $\Omega$ , X L = 20 $\Omega$ , X C = 12 $\Omega$ . For a supply of 100 V, what is the magnitude of the current and its nature?} \hfill [Topic: steady state analysis of ac circuits consisting of RL, RC and RLC (series) | Hard]
\begin{enumerate}[label=(\alph*),topsep=2pt,itemsep=1pt,partopsep=0pt,parsep=0pt]
    \item {5 A, leading}
    \item {10.6 A, lagging}
    \item {20 A, lagging}
    \item {10.6 A, leading}
\end{enumerate}
\vspace{2mm}

\noindent\textbf{Q52. A series RLC circuit resonates at 500 Hz with a bandwidth of 50 Hz. What are the approximate half-power frequencies?} \hfill [Topic: resonance in series RLC circuit | Hard]
\begin{enumerate}[label=(\alph*),topsep=2pt,itemsep=1pt,partopsep=0pt,parsep=0pt]
    \item {475 Hz and 525 Hz}
    \item {495 Hz and 505 Hz}
    \item {450 Hz and 550 Hz}
    \item {400 Hz and 600 Hz}
\end{enumerate}
\vspace{2mm}

\noindent\textbf{Q53. Two voltages are v 1 = 100 sin ( $\omega$ t ) and v 2 = 100 sin ( $\omega$ t - $\pi$ /3 ) . By what time (in terms of period T ) does v 2 lag v 1 ?} \hfill [Topic: phase difference | Hard]
\begin{enumerate}[label=(\alph*),topsep=2pt,itemsep=1pt,partopsep=0pt,parsep=0pt]
    \item {T /12}
    \item {T /3}
    \item {T /6}
    \item {T /2}
\end{enumerate}
\vspace{2mm}

\noindent\textbf{Q54. Two impedances Z 1 = 10$\angle$6 0 $^\circ$ $\Omega$ and Z 2 = 10$\angle$ - 3 0 $^\circ$ $\Omega$ are connected in parallel. What is the magnitude of the equivalent impedance?} \hfill [Topic: complex representation of impedance | Hard]
\begin{enumerate}[label=(\alph*),topsep=2pt,itemsep=1pt,partopsep=0pt,parsep=0pt]
    \item {14.14 $\Omega$}
    \item {5 $\Omega$}
    \item {7.07 $\Omega$}
    \item {10 $\Omega$}
\end{enumerate}
\vspace{2mm}

\noindent\textbf{Q55. A coil draws 10 A when connected to a 120 V DC supply and 10 A when connected to a 230 V, 50 Hz AC supply. What is the power factor of the coil on AC?} \hfill [Topic: power factor and power calculation in RL, RC and RLC circuits | Hard]
\begin{enumerate}[label=(\alph*),topsep=2pt,itemsep=1pt,partopsep=0pt,parsep=0pt]
    \item {0.85}
    \item {0.52}
    \item {0.71}
    \item {0.60}
\end{enumerate}
\vspace{2mm}

\noindent\textbf{Q56. In a series RLC circuit, the applied voltage is 200 V, V R = 120 V and V L = 160 V. If V C = 0 , this is inconsistent; given V C = 320 V instead, what is V R consistent with 200 V supply?} \hfill [Topic: steady state analysis of ac circuits consisting of RL, RC and RLC (series) | Hard]
\begin{enumerate}[label=(\alph*),topsep=2pt,itemsep=1pt,partopsep=0pt,parsep=0pt]
    \item {160 V}
    \item {200 V}
    \item {120 V}
    \item {80 V}
\end{enumerate}
\vspace{2mm}

\noindent\textbf{Q57. A balanced delta-connected load of Z = 15$\angle$3 0 $^\circ$ $\Omega$ per phase is connected to a 400 V line. What is the line current?} \hfill [Topic: relations in line and phase voltages and currents in star and delta networks | Hard]
\begin{enumerate}[label=(\alph*),topsep=2pt,itemsep=1pt,partopsep=0pt,parsep=0pt]
    \item {80 A}
    \item {46.2 A}
    \item {26.7 A}
    \item {15.4 A}
\end{enumerate}
\vspace{2mm}

\noindent\textbf{Q58. A balanced three-phase star load draws a line current of 20 A at 400 V line voltage with power factor 0.866 . What is the total real power consumed?} \hfill [Topic: power factor and power calculation in RL, RC and RLC circuits | Hard]
\begin{enumerate}[label=(\alph*),topsep=2pt,itemsep=1pt,partopsep=0pt,parsep=0pt]
    \item {24 kW}
    \item {8 kW}
    \item {6.93 kW}
    \item {12 kW}
\end{enumerate}
\vspace{2mm}

\noindent\textbf{Q59. At resonance in a series RLC circuit, the current is 2 A and R = 20 $\Omega$ . If L = 50 mH and the resonant frequency is 1000 rad/s, what is the voltage across the capacitor?} \hfill [Topic: resonance in series RLC circuit | Hard]
\begin{enumerate}[label=(\alph*),topsep=2pt,itemsep=1pt,partopsep=0pt,parsep=0pt]
    \item {200 V}
    \item {100 V}
    \item {40 V}
    \item {50 V}
\end{enumerate}
\vspace{2mm}

\noindent\textbf{Q60. A sinusoidal voltage v ( t ) = 170 sin ( 377 t + 3 0 $^\circ$ ) V is applied. What are the amplitude, frequency, and phase respectively?} \hfill [Topic: definitions of amplitude, phase | Hard]
\begin{enumerate}[label=(\alph*),topsep=2pt,itemsep=1pt,partopsep=0pt,parsep=0pt]
    \item {170 V, 377 Hz, 3 0 $^\circ$}
    \item {170 V, 60 Hz, 3 0 $^\circ$}
    \item {120 V, 60 Hz, 3 0 $^\circ$}
    \item {170 V, 50 Hz, 0.52 rad}
\end{enumerate}
\vspace{2mm}


\newpage
\section{Section B: Comprehensive Analytical Problems with Step-by-Step Solutions}

\begin{subjectiveq}{2}{1}
A series $RLC$ circuit consists of resistance $R = 10\,\Omega$, inductance $L = 50\text{ mH}$, and capacitance $C = 20\,\mu\text{F}$ connected across a sinusoidal voltage supply $v(t) = 230\sqrt{2}\sin(100\pi t)\text{ V}$. Determine:\begin{enumerate}[label=(\alph*)]\item The resonant frequency $f_0$, quality factor $Q_0$, and bandwidth $BW$ of the circuit.\item The circuit impedance $\mathbf{Z}$, current $\mathbf{I}$, power factor $\cos\phi$, and real power $P$ absorbed at supply frequency $f = 50\text{ Hz}$.\item The half-power frequencies $f_1$ and $f_2$.\end{enumerate}
\end{subjectiveq}

\begin{solbox}
\textbf{Step 1: Resonant Frequency, Quality Factor, and Bandwidth}\\The resonant angular frequency is:\[ \omega_0 = \frac{1}{\sqrt{LC}} = \frac{1}{\sqrt{50 \times 10^{-3} \times 20 \times 10^{-6}}} = 1000\text{ rad/s} \]The resonant cyclic frequency is $f_0 = \frac{\omega_0}{2\pi} = \frac{1000}{2\pi} \approx 159.15\text{ Hz}$.The Quality Factor $Q_0$ is given by:\[ Q_0 = \frac{\omega_0 L}{R} = \frac{1000 \times 0.050}{10} = 5.0 \]The Bandwidth in Hertz is:\[ BW = \frac{f_0}{Q_0} = \frac{159.15}{5.0} \approx 31.83\text{ Hz} \]\textbf{Step 2: Circuit Behavior at $f = 50\text{ Hz}$ ($\omega = 100\pi \approx 314.16\text{ rad/s}$)}\\Compute inductive and capacitive reactances:\begin{align*}X_L &= \omega L = 314.16 \times 0.050 = 15.71\,\Omega \\X_C &= \frac{1}{\omega C} = \frac{1}{314.16 \times 20 \times 10^{-6}} \approx 159.15\,\Omega\end{align*}Net reactance is $X = X_L - X_C = 15.71 - 159.15 = -143.44\,\Omega$ (capacitive).Complex impedance: $\mathbf{Z} = 10 - j143.44\,\Omega = 143.79 \angle -86.0^\circ\,\Omega$.RMS current magnitude: $I = \frac{V_{rms}}{|\mathbf{Z}|} = \frac{230}{143.79} \approx 1.60\text{ A}$.Power factor: $\cos\phi = \frac{R}{|\mathbf{Z}|} = \frac{10}{143.79} \approx 0.0695\text{ leading}$.Real power: $P = V I \cos\phi = 230 \times 1.60 \times 0.0695 \approx 25.6\text{ W}$ ($P = I^2 R = 1.60^2 \times 10 = 25.6\text{ W}$).\\[2mm]\textbf{Step 3: Half-Power Frequencies}\\\[ f_1 = f_0 - \frac{BW}{2} = 159.15 - 15.92 = 143.23\text{ Hz}, \quad f_2 = f_0 + \frac{BW}{2} = 159.15 + 15.92 = 175.07\text{ Hz} \]
\end{solbox}

\begin{examinertip}
\textbf{Key Examiner Caveat}: At frequencies below resonance ($\omega < \omega_0$), $X_C > X_L$, making the circuit predominantly capacitive and current leading.
\end{examinertip}
\vspace{4mm}
